% Distilled blueprint for Evans, Appendix B (Inequalities).

\section{Convex Functions}
\label{sec:convex-functions}

\begin{definition}[Convex function]
\label{def:convex-function-rn}
\dcref{appB:B.1-def-1}
\group{Convex Analysis}
A function $f:\R^n\to\R$ is \textit{convex} if
\[
(1)\qquad
f(\tau x+(1-\tau)y)
\leq \tau f(x)+(1-\tau)f(y)
\]
for all $x,y\in\R^n$ and $0\leq\tau\leq1$.
\end{definition}

\begin{theorem}[Supporting hyperplanes]
\label{thm:supporting-hyperplanes}
\uses{def:convex-function-rn}
\dcref{appB:B.1-thm-1}
\group{Convex Analysis}

If $f:\R^n\to\R$ is convex, then for every $x\in\R^n$ there exists $r\in\R^n$ such that
\[
(2)\qquad f(y)\geq f(x)+r\cdot(y-x)
\]
for every $y\in\R^n$.
\end{theorem}

\begin{remark}[Supporting hyperplanes and uniform convexity]
\label{rem:supporting-hyperplane-remarks}
\uses{thm:supporting-hyperplanes}
\dcref{appB:B.1-rem-1}
\group{Convex Analysis}

(i) The affine map $y\mapsto f(x)+r\cdot(y-x)$ is a supporting hyperplane to $f$ at $x$, and (2) places the graph of $f$ above it. If $f$ is differentiable at $x$, then $r=Df(x)$.

(ii) For $f\in C^2(\R^n)$, convexity is equivalent to $D^2f\geq0$. The function is \textit{uniformly convex} if $D^2f\geq\theta I$ for some $\theta>0$, equivalently
\[
\sum_{i,j=1}^n f_{x_ix_j}(x)\xi_i\xi_j
\geq\theta|\xi|^2
\qquad(x,\xi\in\R^n).
\]
\end{remark}

\begin{theorem}[Jensen's inequality]
\label{thm:jensens-inequality}
\uses{def:convex-function-rn, thm:supporting-hyperplanes, def:function-notation-basic}
\dcref{appB:B.1-thm-2}
\group{Convex Analysis}

Let $f:\R\to\R$ be convex, let $U\subset\R^n$ be open and bounded, and let $u:U\to\R$ be summable. Then
\[
(3)\qquad
f\left(\fint_Uu\,dx\right)
\leq\fint_Uf(u)\,dx,
\qquad
\fint_Uu\,dx:=\frac1{|U|}\int_Uu\,dx.
\]
\end{theorem}
\begin{proof}
\sketch Set $p=\fint_Uu\,dx$ and choose a supporting slope $r$ at $p$. Then
\[
f(u(x))\geq f(p)+r(u(x)-p).
\]
Average over $U$; the affine error has average zero, giving (3).
\end{proof}

\section{Elementary Inequalities}
\label{sec:elementary-inequalities}

\begin{lemma}[Cauchy's inequality]
\label{lem:cauchys-inequality}
\dcref{appB:B.2-a}
\group{Elementary Inequalities}
For all $a,b\in\R$,
\[
(4)\qquad ab\leq\frac{a^2}{2}+\frac{b^2}{2}.
\]
\end{lemma}
\begin{proof}
\sketch Expand $0\leq(a-b)^2=a^2-2ab+b^2$.
\end{proof}

\begin{lemma}[Cauchy's inequality with $\epsilon$]
\label{lem:cauchys-inequality-epsilon}
\uses{lem:cauchys-inequality}
\dcref{appB:B.2-b}
\group{Elementary Inequalities}

For $a,b>0$ and $\epsilon>0$,
\[
(5)\qquad
ab\leq\epsilon a^2+\frac{b^2}{4\epsilon}.
\]
\end{lemma}
\begin{proof}
\sketch Apply (4) to $(2\epsilon)^{1/2}a$ and $b/(2\epsilon)^{1/2}$.
\end{proof}

\begin{lemma}[Young's inequality]
\label{lem:youngs-inequality}
\uses{def:convex-function-rn}
\dcref{appB:B.2-c}
\group{Elementary Inequalities}

Let $1<p,q<\infty$ and $p^{-1}+q^{-1}=1$. Then for $a,b>0$,
\[
(6)\qquad ab\leq\frac{a^p}{p}+\frac{b^q}{q}.
\]
\end{lemma}
\begin{proof}
\sketch Convexity of the exponential gives
\[
ab=e^{p^{-1}\log a^p+q^{-1}\log b^q}
\leq p^{-1}e^{\log a^p}+q^{-1}e^{\log b^q},
\]
which is (6).
\end{proof}

\begin{lemma}[Young's inequality with $\epsilon$]
\label{lem:youngs-inequality-epsilon}
\uses{lem:youngs-inequality}
\dcref{appB:B.2-d}
\group{Elementary Inequalities}

Let $1<p,q<\infty$ satisfy $p^{-1}+q^{-1}=1$. For $a,b>0$ and $\epsilon>0$,
\[
(7)\qquad
ab\leq\epsilon a^p+C(\epsilon)b^q
\qquad(a,b>0),
\]
where
\[
C(\epsilon)=(\epsilon p)^{-q/p}q^{-1}.
\]
\end{lemma}
\begin{proof}
\sketch Apply (6) to $(\epsilon p)^{1/p}a$ and $b/(\epsilon p)^{1/p}$.
\end{proof}

\begin{lemma}[H\"older's inequality]
\label{lem:holders-inequality}
\uses{lem:youngs-inequality}
\dcref{appB:B.2-e}
\group{Elementary Inequalities}

Let $1\leq p,q\leq\infty$ satisfy $p^{-1}+q^{-1}=1$. If $u\in L^p(U)$ and $v\in L^q(U)$, then
\[
(8)\qquad
\int_U|uv|\,dx
\leq\|u\|_{L^p(U)}\|v\|_{L^q(U)}.
\]
\end{lemma}
\begin{proof}
\sketch For $1<p,q<\infty$, homogeneity reduces to $\|u\|_{L^p}=\|v\|_{L^q}=1$. Apply (6) pointwise to $|u|$ and $|v|$ and integrate. The cases $(p,q)=(1,\infty)$ and $(\infty,1)$ follow directly from the essential-supremum definition.
\end{proof}

\begin{lemma}[Minkowski's inequality]
\label{lem:minkowskis-inequality}
\uses{lem:holders-inequality}
\dcref{appB:B.2-f}
\group{Elementary Inequalities}

If $1\leq p\leq\infty$ and $u,v\in L^p(U)$, then
\[
(9)\qquad
\|u+v\|_{L^p(U)}
\leq\|u\|_{L^p(U)}+\|v\|_{L^p(U)}.
\]
\end{lemma}
\begin{proof}
\sketch For $1\leq p<\infty$, write
\[
\int_U|u+v|^p\,dx
\leq\int_U|u+v|^{p-1}(|u|+|v|)\,dx
\]
and apply H\"older separately to the two terms. This bounds the right side by
\[
\|u+v\|_{L^p}^{p-1}
(\|u\|_{L^p}+\|v\|_{L^p});
\]
division yields (9), with the zero-norm case immediate. For $p=\infty$, use the pointwise triangle inequality and take the essential supremum.
\end{proof}

\begin{remark}[Discrete H\"older and Minkowski inequalities]
\label{rem:discrete-holder-minkowski-inequalities}
\uses{lem:holders-inequality, lem:minkowskis-inequality}
\dcref{appB:B.2-f-rem}
\group{Elementary Inequalities}

For $a,b\in\R^n$, $1\leq p<\infty$, and $p^{-1}+q^{-1}=1$,
\[
(10)\qquad
\begin{cases}
\displaystyle
\left|\sum_{k=1}^na_kb_k\right|
\leq
\left(\sum_{k=1}^n|a_k|^p\right)^{1/p}
\left(\sum_{k=1}^n|b_k|^q\right)^{1/q},\\[1em]
\displaystyle
\left(\sum_{k=1}^n|a_k+b_k|^p\right)^{1/p}
\leq
\left(\sum_{k=1}^n|a_k|^p\right)^{1/p}
+\left(\sum_{k=1}^n|b_k|^p\right)^{1/p}.
\end{cases}
\]
\end{remark}

\begin{lemma}[General H\"older inequality]
\label{lem:general-holder-inequality}
\uses{lem:holders-inequality}
\dcref{appB:B.2-g}
\group{Elementary Inequalities}

Let $1\leq p_1,\ldots,p_m\leq\infty$ satisfy
\[
\sum_{k=1}^m\frac1{p_k}=1,
\]
and suppose $u_k\in L^{p_k}(U)$ for $k=1,\ldots,m$. Then
\[
(11)\qquad
\int_U|u_1\cdots u_m|\,dx
\leq\prod_{k=1}^m\|u_k\|_{L^{p_k}(U)}.
\]
\end{lemma}
\begin{proof}
\sketch Induct on $m$, applying \cref{lem:holders-inequality} at each step with the conjugate exponents determined by the remaining reciprocal sum.
\end{proof}

\begin{lemma}[Interpolation inequality for $L^p$-norms]
\label{lem:interpolation-inequality-lp-norms}
\uses{lem:holders-inequality}
\dcref{appB:B.2-h}
\group{Elementary Inequalities}

Assume $1\leq s\leq r\leq t\leq\infty$ and
\[
\frac1r=\frac\theta s+\frac{1-\theta}{t}.
\]
If $u\in L^s(U)\cap L^t(U)$, then $u\in L^r(U)$ and
\[
(12)\qquad
\|u\|_{L^r(U)}
\leq\|u\|_{L^s(U)}^\theta
\|u\|_{L^t(U)}^{1-\theta}.
\]
\end{lemma}
\begin{proof}
\sketch In the nondegenerate case, factor
\[
|u|^r=|u|^{\theta r}|u|^{(1-\theta)r}
\]
and apply H\"older with conjugate exponents $s/(\theta r)$ and $t/((1-\theta)r)$. Taking the $r$-th root gives (12). If $\theta\in\{0,1\}$ or an endpoint exponent is infinite, the same conclusion follows by omitting the corresponding factor and using its essential supremum.
\end{proof}

\begin{lemma}[Cauchy--Schwarz inequality]
\label{lem:cauchy-schwarz-inequality}
\dcref{appB:B.2-i}
\group{Elementary Inequalities}
For $x,y\in\R^n$,
\[
(13)\qquad |x\cdot y|\leq|x||y|.
\]
\end{lemma}
\begin{proof}
\sketch For $\epsilon>0$,
\[
0\leq|x\pm\epsilon y|^2
=|x|^2\pm2\epsilon x\cdot y+\epsilon^2|y|^2,
\]
so
\[
\pm x\cdot y
\leq\frac{|x|^2}{2\epsilon}+\frac\epsilon2|y|^2.
\]
If $y\ne0$, minimize at $\epsilon=|x|/|y|$; if $y=0$, the conclusion is immediate.
\end{proof}

\begin{remark}[Cauchy--Schwarz inequality for a symmetric nonnegative matrix]
\label{rem:cauchy-schwarz-matrix-form}
\uses{lem:cauchy-schwarz-inequality}
\dcref{appB:B.2-i-rem}
\group{Elementary Inequalities}

If $A=((a_{ij}))$ is a symmetric nonnegative $n\times n$ matrix, then for $x,y\in\R^n$,
\[
(14)\qquad
\left|\sum_{i,j=1}^na_{ij}x_iy_j\right|
\leq
\left(\sum_{i,j=1}^na_{ij}x_ix_j\right)^{1/2}
\left(\sum_{i,j=1}^na_{ij}y_iy_j\right)^{1/2}.
\]
\end{remark}

\begin{lemma}[Gronwall's inequality (differential form)]
\label{lem:gronwalls-inequality-differential-form}
\dcref{appB:B.2-j}
\group{Gronwall's Inequality}
(i) Let $\eta:[0,T]\to[0,\infty)$ be absolutely continuous and suppose
\[
(15)\qquad
\eta'(t)\leq\phi(t)\eta(t)+\psi(t)
\]
for almost every $t\in[0,T]$, where $\phi$ and $\psi$ are nonnegative summable functions. Then
\[
(16)\qquad
\eta(t)
\leq e^{\int_0^t\phi(s)\,ds}
\left(\eta(0)+\int_0^t\psi(s)\,ds\right)
\]
for every $0\leq t\leq T$.

(ii) In particular, if $\eta'\leq\phi\eta$ on $[0,T]$ and $\eta(0)=0$, then $\eta\equiv0$ on $[0,T]$.
\end{lemma}
\begin{proof}
\sketch Multiply (15) by the integrating factor $e^{-\int_0^t\phi}$. For almost every $s$,
\[
\frac d{ds}\left(\eta(s)e^{-\int_0^s\phi(r)\,dr}\right)
\leq e^{-\int_0^s\phi(r)\,dr}\psi(s).
\]
Integrate from $0$ to $t$ and use $e^{-\int_0^s\phi}\leq1$ to obtain (16). Setting $\psi=0$ and $\eta(0)=0$ gives (ii).
\end{proof}

\begin{lemma}[Gronwall's inequality (integral form)]
\label{lem:gronwalls-inequality-integral-form}
\uses{lem:gronwalls-inequality-differential-form}
\dcref{appB:B.2-k}
\group{Gronwall's Inequality}

(i) Let $\xi:[0,T]\to[0,\infty)$ be summable and suppose
\[
(17)\qquad
\xi(t)\leq C_1\int_0^t\xi(s)\,ds+C_2
\]
for almost every $t\in[0,T]$, where $C_1,C_2\geq0$. Then
\[
(18)\qquad
\xi(t)\leq C_2(1+C_1te^{C_1t})
\]
for almost every $0\leq t\leq T$.

(ii) In particular, if
\[
\xi(t)\leq C_1\int_0^t\xi(s)\,ds
\]
for almost every $t\in[0,T]$, then $\xi(t)=0$ almost everywhere.
\end{lemma}
\begin{proof}
\sketch Set $\eta(t)=\int_0^t\xi(s)\,ds$. Then $\eta(0)=0$ and (17) gives
\[
\eta'(t)\leq C_1\eta(t)+C_2
\]
almost everywhere. Differential Gronwall yields $\eta(t)\leq C_2te^{C_1t}$; substitution into (17) gives (18). Taking $C_2=0$ proves (ii).
\end{proof}
